\subsection{The complete ninth-root class group and cyclotomic capitulation}

\begin{statement}{A.206}{Proved}
{An exact regulator index and the four-dimensional three-class group}
Retain the actual fields and automorphisms of A.191 and A.203:
\[
\begin{gathered}
 F=\Q(\alpha,\zeta),\quad F'=\Q(\alpha,z),\quad
 E_2=\Q(t,\zeta),\quad M=\Q(t,z),\\
 \alpha^3=3,\quad t^3=d^2\alpha,\quad d=649,
 \quad z=\zeta_9,\quad\zeta=z^3,\\
 \sigma(t)=\zeta t,\quad\sigma(z)=z,
 \qquad h(t)=t,\quad h(z)=z^4.
\end{gathered}
\tag{A.206.1}
\]
Put \(A=\operatorname{Cl}(M)\{3\}\),
\(A_2=\operatorname{Cl}(E_2)\{3\}\), and
\(\pi=\sigma-1\).  The exact class group and its
\(\mathcal O=\Z_3[\zeta]\)-module structure are
\[
 A\simeq\mathcal O/(\pi^2)\oplus
             \mathcal O/(\pi)\oplus\mathcal O/(\pi),
 \qquad A\simeq(\Z/3)^4
          \quad\hbox{as an abelian group}.
\tag{A.206.2}
\]
Here \(\zeta\) acts through the actual automorphism
\(\sigma\).  Ideal norm to \(E_2\) induces the
isomorphism
\[
 \overline N_h:A/(\sigma-1)A
            \xrightarrow{\ \sim\ }A_2\simeq(\Z/3)^3.
\tag{A.206.3}
\]
In particular \(N_h:A\to A_2\) is onto and its kernel
is the order-three group \((\sigma-1)A\).

The actual ideal extension map
\(j_h:A_2\to A\) is zero: every three-primary ideal
class of \(E_2\) capitulates in \(M\).
The action of \(h-1\) on \(A\) has rank one over
\(\F_3\) and square zero.  These assertions use the
literal ideal extension and residue-weighted ideal norm,
whose composites are
\[
 N_hj_h=3\quad\hbox{on }A_2,\qquad
 j_hN_h=1+h+h^2\quad\hbox{on }A.
\tag{A.206.4}
\]
The subgroup and quotient maps in (A.206.3)--(A.206.4)
are retained throughout the class-number calculation below.
Let \(\mathfrak P_3\) be the unique prime of \(M\)
over three, and write \(c_3\) for the three-primary
projection of its ideal class.  Then \(c_3\) has exact
order three and
\[
 \ker N_h=(\sigma-1)A=(h-1)A=\langle c_3\rangle.
\tag{A.206.5}
\]
For \(S\) the primes above \(2,3,11,59\) and the
archimedean places, the norm induces the exact isomorphism
\[
 \operatorname{Cl}_S(M)\{3\}
       =A/\langle c_3\rangle
       \xrightarrow{\ \sim\ }A_2\simeq(\Z/3)^3.
\tag{A.206.6}
\]
\end{statement}

\begin{proof}
\emph{The four embedded cubic subfields and their units.}
The two commuting generators in (A.206.1) give
\(G=\operatorname{Gal}(M/F)=\langle\sigma,h\rangle
\simeq(\Z/3)^2\).  Let \(g(t)=zt\), \(g(z)=z\).
It is an actual automorphism of \(M\), with
\(g^3=\sigma\).  The four order-three subgroup fixed
fields of \(G\) are
\[
 L_0=F',\qquad L_{j+1}=g^jE_2
       =\Q(tz^j,\zeta)\quad(0\leq j\leq2).
\tag{A.206.7}
\]
More explicitly \(L_0=M^{\langle\sigma\rangle}\),
and \(L_{j+1}=M^{\langle\sigma^{-j}h\rangle}\),
since \((\sigma^{-j}h)(tz^j)=tz^j\).
The automorphisms \(g^j\) identify the last three
fields, their units and their ideal class groups.

All these fields are totally imaginary.  Their numbers
of complex pairs, free unit ranks, and numbers of roots
of unity are
\[
\begin{array}{c|c|c|c}
 L&r_2(L)&r_2(L)-1&w_L\\\hline
 F&3&2&6\\
 F'&9&8&18\\
 L_1,L_2,L_3&9&8&6\\
 M&27&26&18
\end{array}.
\tag{A.206.8}
\]
For the last column, the full holomorph of A.191 makes
\(\Q(z)\) the maximal abelian subfield of \(M/\Q\).
Thus every root of unity in \(M\) belongs to
\(\Q(z)\), whose roots of unity form \(\mu_{18}\).
The intersection of \(L_j\), \(j\geq1\), with
\(\Q(z)\) is \(\Q(\zeta)\): adjoining \(z\)
to \(L_j\) is the actual cubic extension in (A.206.7).
This proves all the asserted root-of-unity orders.
Only total imaginaryness is used in the logarithmic argument.

Let \(B_{\rm act}\) be the subgroup of
\(\mathcal O_M^\times\) generated by the four actual
\(\mathcal O_{L_i}^\times\).  The literal group-ring
identity of A.194 is
\[
 \sum_{i=0}^3\sum_{a\in\operatorname{Gal}(M/L_i)}a
                 =3+\sum_{a\in G}a.
\tag{A.206.9}
\]
Applied multiplicatively to any actual unit, it shows
that its cube lies in \(B_{\rm act}\); the full norm
on the right is an \(F\)-unit and already belongs
to this subgroup.  Thus the actual quotient is finite
and killed by three.  Passing to pro-three completions
does not change it.  A.199's complete residue detector
therefore gives the exact ordinary index
\[
       [\mathcal O_M^\times:B_{\rm act}]=3^6.
\tag{A.206.10}
\]
This is an index of the full subfield unit groups.
The roots of unity of \(M\) already belong to \(F'\),
so the same index occurs in their free logarithmic lattices.

\emph{The intersection of each log lattice with the base space.}
For a totally imaginary field \(L\), use the logarithmic
map
\(u\mapsto(2\log|\iota(u)|)_\iota\), with one coordinate
for each complex pair, in its sum-zero Euclidean hyperplane.
Write \(\Lambda_L\) for the image of its actual global
units.  Embed all the log spaces into the log space of
\(M\) by repeating coordinates along field embeddings.
If a unit \(u\in\mathcal O_{L_i}^\times\) has logarithm
in the base real log space \(V_F\), then for a generator
\(a\) of \(\operatorname{Gal}(L_i/F)\), every conjugate
of \(a(u)/u\) has absolute value one.  It is an algebraic
integer unit, so Kronecker's theorem makes it a root of unity.
Indeed the monic polynomials of all its powers have bounded
degree and bounded integer coefficients; only finitely many
roots occur, so two powers agree.

For \(L_i=L_{j+1}\), its norm-one condition puts this
ratio in \(\mu_3\), and \(u^3\) is an actual \(F\)-unit.
The Kummer restriction kernel for \(F\to L_{j+1}\) is
the line generated by
\((tz^j)^3=d^2\alpha\zeta^j\).  At each prime of
\(F\) above eleven this radicand has valuation two.
Neither of its nonzero powers modulo cubes can therefore
be a unit class.  Hence \(u^3=v^3\) for a unit \(v\)
of \(F\); then \(u/v\in\mu_3\subset F\).

For \(L_0=F'\), write \(h(u)/u=z^k\).  Its possible
sign is positive because the norm is one and the degree
is odd.  The literal norm is
\[
        N_{F'/F}(z^k)=z^{(1+4+7)k}=\zeta^k.
\tag{A.206.11}
\]
Thus \(3\mid k\).  Dividing \(u\) by a suitable
power of \(z\) makes it \(h\)-fixed, without changing
its logarithm.  These arguments prove the exact intersection
in every one of the four embedded log spaces:
\[
                   \Lambda_{L_i}\cap V_F=\Lambda_F.
\tag{A.206.12}
\]

\emph{The complete regulator factor.}
Orthogonal character decomposition for the actual permutation
action of \(G\) on the 27 complex pairs gives
\[
 V_M=V_F\mathbin\perp W_0\mathbin\perp W_1
                   \mathbin\perp W_2\mathbin\perp W_3,
 \quad \dim V_F=2,\quad\dim W_i=6,
 \quad V_{L_i}=V_F\mathbin\perp W_i.
\tag{A.206.13}
\]
The summation map from \(\bigoplus_i V_{L_i}\) to
\(V_M\) has lattice kernel exactly the four tuples
of elements of \(\Lambda_F\) whose sum is zero, by
(A.206.12) and the distinct orthogonal \(W_i\).
This kernel is \(\Lambda_F\otimes A_3\), where the
root lattice \(A_3\) has Gram determinant four.
Its covolume is therefore
\(4\operatorname{covol}(\Lambda_F)^3\).
On the orthogonal complement of the kernel, the sum
map multiplies lengths by two on the two-dimensional
base summand and is an isometry on every \(W_i\).
Its absolute determinant is four.  These two factors
cancel in the exact lattice quotient formula.  Consequently
\[
 \operatorname{covol}(\log B_{\rm act})
   =\frac{\prod_i\operatorname{covol}_M(\Lambda_{L_i})}
          {\operatorname{covol}_M(\Lambda_F)^3}.
\tag{A.206.14}
\]
Here the subscript indicates the metric induced from \(M\).

The embedding \(L_i\to M\) repeats every log coordinate
three times; its rank-eight covolume is multiplied by
\((\sqrt3)^8=3^4\).  The embedding \(F\to M\)
repeats each coordinate nine times; its rank-two covolume
is multiplied by \(3^2\).  In native metrics (A.206.14)
thus contributes the exact factor \(3^{16-6}=3^{10}\).
The native log covolume is \(\sqrt{r_2(L)}R_L\), where
\(R_L\) is the usual regulator: the degree-zero lattice
in \(\Z^{r_2(L)}\) has covolume \(\sqrt{r_2(L)}\).
Using (A.206.8), the remaining square-root factor is
\[
       \frac{(\sqrt9)^4}{\sqrt{27}(\sqrt3)^3}=3.
\tag{A.206.15}
\]
Dividing (A.206.14) by the full log-lattice covolume
and using (A.206.10) proves
\[
 \frac{\prod_{i=0}^3 R_{L_i}}{R_M R_F^3}
  =\frac{[\mathcal O_M^\times:B_{\rm act}]}{3^{11}}
  =3^{-5}.
\tag{A.206.16}
\]
All base intersections, metric changes and torsion kernels
have been included in this equality.

\emph{The zeta identity and the class group.}
The actual permutation characters satisfy
\(\Q[G]+3\Q=\bigoplus_i\Q[G/H_i]\), where
\(H_i=\operatorname{Gal}(M/L_i)\).
Taking their Euler factors, including inertia invariants
at ramified primes, gives
\[
                \zeta_M(s)\zeta_F(s)^3
                     =\prod_{i=0}^3\zeta_{L_i}(s).
\tag{A.206.17}
\]
The analytic class-number formula at zero, in
\cite[Section 2, Theorem 2.1, printed p.~304]{BIR_Darmon1995},
states that the leading coefficient is \(-h_LR_L/w_L\)
at its zero of order \(r_1(L)+r_2(L)-1\).
In that source take \(S\) to be exactly the archimedean
places and \(T\) empty; its degree-zero log-lattice
regulator is precisely the one used above.
The zero orders in (A.206.17) agree:
\(26+3\cdot2=4\cdot8\).  The four leading signs
also agree.  The root-of-unity ratio is exactly
\(18\cdot6^3/(18\cdot6^3)=1\).
It follows that
\[
       \frac{h_Mh_F^3}{\prod_i h_{L_i}}
              =\frac{\prod_iR_{L_i}}{R_MR_F^3}=3^{-5}.
\tag{A.206.18}
\]
A.138 and A.191 give \(\operatorname{Cl}(F)\{3\}=0\)
and \(\operatorname{Cl}(F')\{3\}=0\).
A.173 gives \(|\operatorname{Cl}(L_i)\{3\}|=27\)
for \(i=1,2,3\), transported by the actual \(g^j\).
Thus (A.206.18) yields
\[
                          |A|=3^{9-5}=3^4.
\tag{A.206.19}
\]

The class norm to \(F'\) has zero three-primary target.
Its composite with ideal extension is
\(1+\sigma+\sigma^2\); therefore this operator kills
\(A\).  It is a finite module over the discrete valuation
ring \(\mathcal O=\Z_3[\zeta]\), with uniformizer
\(\pi=\sigma-1\).  A.191 proves
\(|A/\pi A|=27\).  The elementary-divisor decomposition
over \(\mathcal O\) consequently has three nonzero
summands whose lengths sum to four by (A.206.19).
Their lengths are exactly \(2,1,1\).
Since \((\pi^2)=(3)\), their additive group is
\((\Z/3)^4\), proving (A.206.2).
The A.191 quotient isomorphism is induced by the other
ideal norm, \(N_h=N_{M/E_2}\), and gives exactly
(A.206.3).

\emph{The actual cyclotomic capitulation map.}
Only the unique three-prime ramifies in \(M/E_2\).
A.203 proves that every actual \(E_2\) global unit
is a field norm.  The same cyclic ambiguous-class
formula used in A.173 therefore gives
\[
       |A^h|=|A_2|\,
          \frac{3}{3[\mathcal O_{E_2}^\times:
                \mathcal O_{E_2}^\times\cap N_hM^\times]}
                    =27.
\tag{A.206.20}
\]
Hence \(h-1\) has rank one on the four-dimensional
\(\F_3\)-space \(A\).  It is nilpotent since
\((h-1)^3=h^3-1=0\); a rank-one nilpotent operator
has square zero.  In characteristic three,
\(1+h+h^2=(h-1)^2=0\).
The ordinary ideal identities (A.206.4) hold before
taking primary parts.  Thus \(j_hN_h=0\), and the
surjectivity in (A.206.3) forces \(j_h=0\).
The other composite is multiplication by three on
the exponent-three target, as (A.206.4) states.

\emph{The unique ramified-prime class and localization.}
For the cyclic group \(\langle h\rangle\), A.203
gives
\(\widehat H^0(\langle h\rangle,\mathcal O_M^\times)
 \simeq(\Z/3)^2\).
The unit Herbrand quotient is \(1/3\), by the same
logarithmic representation calculation used in A.173:
the rational unit representation is
\(\Q[\langle h\rangle]^9-\Q\), and roots of unity
have Herbrand quotient one.  It follows that
\[
       |H^1(\langle h\rangle,\mathcal O_M^\times)|=27.
\tag{A.206.21}
\]
The usual capitulation map is literal here.  An ideal
\(\mathfrak a\) of \(E_2\) whose extension is
\((w)\) gives the unit cocycle \(h(w)/w\).
Changing its principal generator changes this by a
unit coboundary.  If that cocycle is a unit coboundary,
then \(w\) divided by a unit belongs to \(E_2\),
so \(\mathfrak a\) is principal.  Thus the actual
capitulation kernel injects into the group in (A.206.21).
Its three-primary subgroup is all of \(A_2\), by the
preceding paragraph, and already has order 27.  It
therefore exhausts that cohomology group.

Hilbert--90 supplies an actual \(w\in M^\times\)
for every such unit cocycle.  The map
\[
 H^1(\langle h\rangle,\mathcal O_M^\times)
       \longrightarrow\Z/3,\qquad
 [h(w)/w]\longmapsto v_{\mathfrak P_3}(w)\pmod3
\tag{A.206.22}
\]
is well-defined: multiplying \(w\) by an \(E_2\)
element changes its valuation by a multiple of the
ramification index three, and changing the unit cocycle
by a unit coboundary changes it by zero.  Every class
from an extended ideal has valuation divisible by three.
The exhaustion just proved makes (A.206.22) the zero map.

If \(c_3\) were zero, the ordinary class of
\(\mathfrak P_3\) would have some finite order \(a\)
prime to three.  An actual generator
\((w)=\mathfrak P_3^a\) would give the unit cocycle
\(h(w)/w\) with nonzero value \(a\bmod3\) in
(A.206.22), a contradiction.  Thus \(c_3\ne0\);
its order is three by (A.206.2).
Its ideal norm is the old three-prime, whose
three-primary class is zero by A.173's three detecting
characters, as explicitly used in A.191.
The kernel in (A.206.3) has order three, so it equals
\(\langle c_3\rangle\).
Both augmentation images lie in this norm kernel.
Their ranks are one, proving (A.206.5).

Finally A.191 proves through the actual residue-weighted
ideal norm that every \(S\)-prime class lies in
\(\ker N_h\).  Since \(S\) contains \(\mathfrak P_3\),
their three-primary subgroup is exactly
\(\langle c_3\rangle\).  This proves (A.206.6).
This argument retains the split two-primes as well:
an individual prime above a split two-prime is not
replaced by the extension of its lower prime.

The retained certificate
\path{code/two_inert_ninth_root_class_certificate.py}
checks the four literal subgroup stabilizers, the
permutation-character and group-ring identities, the
exact integral log-lattice index \(3^{11}\) in the
native degree-zero lattice model, all regulator and
torsion factors, and the finite module/transfer
identities used above.  The source uses A.199's full
unit index, not a bounded search for class relations.
\end{proof}
